\documentclass[conference]{IEEEtran}
\usepackage[utf8]{inputenc}
\usepackage[T1]{fontenc}
\usepackage[spanish,es-tabla]{babel}
\usepackage{url}
\usepackage{booktabs}
\usepackage{amsmath}
\usepackage[hidelinks]{hyperref}
\usepackage{etoolbox}

% Sustituir los separadores largos que IEEEtran agrega automáticamente.
\patchcmd{\abstract}{---}{\space-\space}{}{%
  \PackageError{paper}{No se pudo ajustar el separador del resumen}{}
}
\patchcmd{\IEEEkeywords}{---}{\space-\space}{}{%
  \PackageError{paper}{No se pudo ajustar el separador de palabras clave}{}
}

% Tres columnas y dos filas; docente primero y después cinco estudiantes.
\newcommand{\authorentry}[3]{%
\begin{minipage}[t]{0.32\textwidth}
\centering\normalfont\small
\textbf{#1}\\
\textit{#2}\\
\textit{Universidad Centroamericana\\José Simeón Cañas}\\
Antiguo Cuscatlán, El Salvador\\[0.3em]
\href{mailto:#3}{#3}
\end{minipage}%
}
\title{Mini-simulador de memoria compartida distribuida a nivel de página en Linux}
\author{%
\begin{tabular}{@{}c@{\hspace{0.02\textwidth}}c@{\hspace{0.02\textwidth}}c@{}}
\authorentry{David Garcia Cerna}
  {Catedrático del Departamento de\\Electrónica e Informática}
  {dgarciac@uca.edu.sv}
&
\authorentry{Armando Antonio\\Aguilar Rosales}
  {Ingeniería Informática}
  {00174323@uca.edu.sv}
&
\authorentry{Josue Daniel\\Rivera Alvarenga}
  {Ingeniería Informática}
  {00048722@uca.edu.sv}
\\[1.5em]
\authorentry{Guillermo Alejandro\\Hernandez Dubon}
  {Ingeniería Informática}
  {00106423@uca.edu.sv}
&
\authorentry{Heber Josue\\Canales Henriquez}
  {Ingeniería Informática}
  {00054621@uca.edu.sv}
&
\authorentry{Mauricio Alejandro\\Contreras Montoya}
  {Ingeniería Informática}
  {00000422@uca.edu.sv}
\end{tabular}}
\begin{document}
\maketitle
\begin{abstract}
Se presenta un simulador educativo de memoria compartida distribuida, construido en
C11 para Linux. Dos procesos independientes se comunican mediante una conexión TCP
local y representan un espacio lógico de 1 MiB dividido en 256 páginas de 4096 bytes.
El sistema detecta accesos a copias inválidas, transfiere páginas completas, migra
propiedad y mantiene copias de lectura opcionales. Un coordinador central serializa
las operaciones y exige invalidaciones confirmadas antes de conceder escritura.
La implementación conserva módulos separados de nodo, transporte, memoria, protocolo
e interfaz. Se verificó con tres programas de pruebas C y 22 escenarios de integración;
nueve escenarios DSM también se ejecutaron con detectores de errores de memoria y
comportamiento indefinido. Se incluyen trazas reproducibles y una demostración de
migración e invalidación. Los fallos de página son simulados, y el sistema no ofrece
recuperación tras caídas ni acceso transparente mediante punteros.
\end{abstract}
\begin{IEEEkeywords}
DSM, coherencia, consistencia secuencial, páginas, TCP, Linux.
\end{IEEEkeywords}

\section{Introducción}
En un sistema distribuido, cada proceso dispone de memoria local y el intercambio de
datos requiere comunicación. Este proyecto permite observar cómo una abstracción de
memoria compartida puede construirse sobre mensajes sin que los participantes compartan
físicamente el mismo arreglo. El objetivo es didáctico: hacer visibles las decisiones
que permiten usar una página remota y modificarla con exclusividad.

La guía del Proyecto 10 solicita dos procesos conectados por sockets, un arreglo
paginado, propiedad inicial concentrada en un nodo y una demo de migración e
invalidación. Se permite ejecutarlos en la misma computadora. La solución utiliza
Pop!\_OS y conserva una arquitectura modular desarrollada previamente para un saludo TCP.

Las contribuciones del trabajo son un simulador ejecutable, un contrato de mensajes
acotado, una CLI explícita y pruebas repetibles con evidencia. No se propone un nuevo
algoritmo de coherencia ni se evalúa rendimiento de sistemas DSM de producción.

\section{Fundamento conceptual}
Li y Hudak estudiaron memoria virtual compartida organizada por páginas y gestores de
coherencia centrales y distribuidos \cite{li}. Nuestro modelo toma como antecedente la
necesidad de mantener información sobre propiedad y copias, pero emplea comandos de
acceso explícito. No implementa el sistema IVY ni acceso transparente a memoria virtual.

La coherencia de página y el orden de accesos son cuestiones distintas. La noción de
consistencia secuencial de Lamport exige explicar los resultados por un orden global
compatible con el orden de cada participante \cite{lamport}. En nuestro diseño, la
serialización del coordinador facilita razonar sobre operaciones completadas sin fallos.

TCP se utiliza como flujo confiable de bytes; no conserva por sí solo las fronteras
de los mensajes de aplicación \cite{tcp}. Por ello el transporte define una cabecera,
acumula recepciones parciales y conserva offsets de envío. Los tratamientos de EOF,
EAGAIN y SIGPIPE siguen las APIs documentadas por Linux \cite{recv,send,poll}.

\section{Modelo y arquitectura}
\subsection{Nodos y recursos}
El mismo binario se ejecuta con identificador 0 o 1. Ambos nodos crean un servidor en
127.0.0.1; únicamente N0 inicia la conexión y N1 la acepta. Esto evita conexiones
TCP duplicadas. El arranque tolera cualquiera de los dos órdenes mediante hasta 30
intentos de conexión, separados por un segundo cuando el peer aún no escucha.

Cada nodo reserva con calloc un arreglo contiguo de $256\times4096=1048576$ bytes.
El número y tamaño de páginas son parámetros del modelo, no páginas físicas consultadas
al kernel. La página $p$ comienza en el desplazamiento $p\times4096$. Se valida el
índice antes de acceder. N0 comienza como dueño de todas las páginas, inicializadas a
cero; N1 mantiene todos sus metadatos en INVALID.

\subsection{Módulos}
Main valida configuración; node administra conexión, eventos y ciclo de vida;
network codifica y transporta bytes; memory aplica operaciones locales; protocol
coordina permisos y transacciones; cli analiza órdenes y presenta resultados.
Memory no conoce sockets y la CLI no modifica directamente los bytes de una página.

Un solo hilo atiende red y entrada estándar mediante poll. La recepción incremental
no espera de forma bloqueante el resto de una trama. Una cola de envío de 65536 bytes
conserva el progreso cuando el socket no acepta más datos. Se limita cada lote de
despacho a 32 mensajes para mantener capacidad de respuesta de la terminal.

\subsection{Metadatos y estados}
Cada página mantiene estado, versión y una indicación informativa del propietario.
Los estados son INVALID, READ\_ONLY y OWNER. Los bytes residuales de INVALID nunca
satisfacen una lectura. OWNER identifica la copia autoritativa; no concede permiso
permanente de escritura si existen lectores.

El directorio autoritativo reside en N0 y contiene propietario, versión confirmada y
máscara de holders válidos. La máscara utiliza un bit por nodo. Las versiones aumentan
una vez por escritura y no cambian al copiar o migrar. Se rechaza una escritura que
desbordaría uint32. Las indicaciones locales de propietario pueden estar desactualizadas
y no se usan como autoridad.

\section{Protocolo implementado}
\subsection{Negociación y mensajes}
Se conserva el saludo textual HELLO/HELLO\_ACK y luego se negocian identidad y geometría
mediante HELLO\_CONFIG. La CLI se habilita tras esta negociación. Cada trama binaria
contiene una cabecera de 40 bytes: diez enteros de 32 bits en orden de red, con magic,
versión, tipo, origen, solicitud, página, propietario, versión de página, flags y longitud.
No se transmiten estructuras C ni su padding.

Los mensajes principales son READ\_REQUEST, WRITE\_REQUEST, FETCH\_PAGE, PAGE\_DATA,
GRANT, INVALIDATE, ACK y ERROR. PAGE\_DATA lleva exactamente 4096 bytes. HELLO\_CONFIG
lleva 12; los restantes no llevan payload. Se comprueba la cabecera antes del cuerpo,
con tipos y longitudes acotados. Un origen y un identificador creciente correlacionan
cada solicitud; no se reutilizan identificadores después de un rechazo BUSY.

\subsection{Lectura con migración y copia opcional}
La orden read cumple el flujo de la guía: si el solicitante necesita datos, se preserva
una instantánea de la página, se invalidan los otros holders y se instala OWNER en el
solicitante. Si el dueño es N1, N0 recupera primero los bytes con FETCH\_PAGE. La
instantánea no se registra como una copia de lectura del coordinador.

La orden separada read-copy mantiene al dueño y concede READ\_ONLY al solicitante.
Esta extensión permite demostrar el escenario condicional de lectores de la guía.
No se presenta una réplica como si fuera migración de propiedad. Una lectura sobre
copia válida evita transferir datos, pero continúa solicitando autorización a N0.

\subsection{Escritura y exclusividad}
Una escritura obtiene primero los datos necesarios y después invalida a todos los
holders distintos del solicitante. N0 envía INVALIDATE a N1 y espera ACK\_INVALIDATED
cuando la copia es remota; si la copia está en N0, se invalida directamente antes de
conceder. Cuando N1 ya es dueño, envía INVALIDATE como solicitud de exclusividad al
coordinador, que invalida su copia local antes de responder GRANT.

El solicitante modifica únicamente después de la concesión. El texto reemplaza toda
la página: se borra el contenido anterior, se copian de 1 a 4095 bytes y se conserva
terminación NUL. ACK\_DONE comunica finalización y versión al coordinador, que actualiza
el directorio. El estado estable vuelve a contener un único dueño.

\subsection{Transacciones y orden}
La máquina global pasa por IDLE, FETCHING si hace falta, INVALIDATING y WAIT\_DONE.
No admite otra operación hasta recibir la finalización. Una solicitud adicional
recibe BUSY y debe reintentarse manualmente. Las solicitudes de N0 recorren la misma
lógica mediante despacho local, evitando reglas distintas para cada nodo.

Los invariantes estables son: un único OWNER por página, coincidencia de datos y versión
entre copias válidas, escritura exclusiva y ausencia de lecturas INVALID. Durante una
transferencia puede haber cero dueños locales; la instantánea y el bloqueo global
impiden perder datos o servir el estado intermedio. No se afirma recuperación de esos
invariantes después de una caída.

\section{Errores y ciclo de vida}
Los argumentos desconocidos, repetidos o incompletos se rechazan. Los puertos deben
estar entre 1 y 65535 y ser distintos. El parser de comandos comprueba rango, longitud
y tokens sobrantes; una orden inválida no modifica memoria. Los datos impresos se
acotan y escapan para evitar interpretar bytes de control como salida de terminal.

El saludo y configuración tienen plazo absoluto de 10 segundos; una operación, trama
parcial o cola pendiente dispone de 5 segundos. Estos plazos no se prolongan por
recibir fragmentos aislados. EOF en reposo cierra normalmente, mientras EOF parcial,
ACK incoherente o timeout abortan la sesión. MSG\_NOSIGNAL permite manejar EPIPE como
error ordinario. Enter, exit y EOF de entrada esperan la operación pendiente; señales
de cierre solicitan aborto y liberación de recursos.

\section{Evaluación reproducible}
\subsection{Entorno y metodología}
Las pruebas se realizaron el 17 de octubre de 2026 en Pop!\_OS 24.04 LTS, Linux
7.1.5 y GCC 13.3, con C11 y advertencias Wall, Wextra y Wpedantic. Se compiló desde
limpio, sin warnings. Las comprobaciones se separan en tres ejecutables C y pruebas
Python que controlan procesos C reales sobre TCP. Python no sustituye el simulador.

Los tests C verifican límites y permisos de memoria, parser, framing y colas. La
integración incluye ambos órdenes de arranque, migraciones, copias, datos actualizados,
longitudes máximas y errores del peer. Un peer controlado retrasa ACK para comprobar
que no se escribe antes de confirmar invalidaciones y provocar BUSY de forma
determinista. Un oráculo secuencial contrasta 48 operaciones con semilla 10.

\begin{table}[t]
\caption{Resultados de verificación}
\centering
\begin{tabular}{lr}
\toprule
Comprobación & Resultado\\
\midrule
Compilación limpia & Sin warnings\\
Programas de prueba C & 3 aprobados\\
Integración completa & 22 aprobados\\
Escenarios DSM con ASan/UBSan & 9 aprobados\\
Demo automática & Aprobada\\
\bottomrule
\end{tabular}
\end{table}

La suite completa tomó 48.166 segundos en la ejecución conservada. Esta duración
incluye esperas intencionales de timeout y reintento y no representa latencia DSM.
AddressSanitizer y UndefinedBehaviorSanitizer no detectaron errores en las ejecuciones
indicadas. Esto es evidencia experimental y no una demostración formal exhaustiva.

\subsection{Demostración y datos observados}
La demo realiza cinco accesos sobre página 3: escritura N0, lectura migratoria N1,
copia de lectura N0, escritura N1 e invalidación de N0, y lectura migratoria de vuelta
a N0. El texto final es el escrito por N1; la versión pasa de 0 a 1 y luego a 2.
Se conservan logs de ambos nodos, comandos, entorno y SHA-256 del código ejecutado.

\begin{table}[t]
\caption{Contadores al terminar la demo}
\centering
\begin{tabular}{lrr}
\toprule
Métrica & N0 & N1\\
\midrule
Fallos simulados & 2 & 1\\
PAGE\_DATA enviados & 1 & 2\\
PAGE\_DATA recibidos & 2 & 1\\
Invalidaciones aplicadas & 2 & 1\\
Commits del coordinador & 5 & -\\
\bottomrule
\end{tabular}
\end{table}

Se transfieren tres páginas completas, 12288 bytes de payload de datos, además de
cabeceras y control. La segunda escritura no retransmite la página porque N1 ya es
dueño. Los contadores de envío representan mensajes encolados; la demo comprueba
su recepción, pero no deben interpretarse como confirmación ante fallos.

\section{Limitaciones y conclusiones}
El sistema cubre los flujos exigidos: dos procesos, detección de ausencia de página,
transferencia completa, instalación, migración y exclusividad con invalidación.
La separación modular y las trazas permiten relacionar cada comando con cambios
observables de estado y mensajes de control.

Se limita a dos nodos en loopback, memoria volátil y una operación activa. No hay
acceso transparente, persistencia, autenticación, recuperación ni reconexión. Tampoco
se mide rendimiento sobre LAN ni se exploran todos los interleavings posibles. El
coordinador es un punto único de fallo. Como trabajo futuro se proponen exploración
de estados, experimentos de red y recuperación diseñada explícitamente, sin asumir
que una reconexión simple preservaría la propiedad.

\begin{thebibliography}{7}
\bibitem{li} K. Li and P. Hudak, ``Memory Coherence in Shared Virtual Memory Systems,''
\emph{ACM Trans. Comput. Syst.}, vol. 7, no. 4, pp. 321-359, Nov. 1989,
doi: 10.1145/75104.75105.
\bibitem{lamport} L. Lamport, ``How to Make a Multiprocessor Computer That Correctly
Executes Multiprocess Programs,'' \emph{IEEE Trans. Comput.}, vol. C-28, no. 9,
pp. 690-691, Sept. 1979.
\bibitem{tcp} Linux man-pages, ``tcp(7).'' [Online]. Available:
\url{https://man7.org/linux/man-pages/man7/tcp.7.html}. Accessed: Oct. 17, 2026.
\bibitem{recv} Linux man-pages, ``recv(2).'' [Online]. Available:
\url{https://man7.org/linux/man-pages/man2/recv.2.html}. Accessed: Oct. 17, 2026.
\bibitem{send} Linux man-pages, ``send(2).'' [Online]. Available:
\url{https://man7.org/linux/man-pages/man2/send.2.html}. Accessed: Oct. 17, 2026.
\bibitem{poll} Linux man-pages, ``poll(2).'' [Online]. Available:
\url{https://man7.org/linux/man-pages/man2/poll.2.html}. Accessed: Oct. 17, 2026.
\end{thebibliography}
\end{document}
